% Sección: Parte 3 – Implementación de PINNs

% Frame 1: Configuración del entorno
\begin{frame}{Configuración del entorno}
  \begin{itemize}
    \item Requisitos previos:
      \begin{itemize}
        \item Python 3.7+ (recomendado 3.8 o 3.9)
        \item Virtualenv o Conda para entornos aislados
        \item CUDA (opcional, para entrenamiento GPU)
      \end{itemize}
    \item Crear entorno con conda:
      \begin{block}{}
        \begin{verbatim}
conda create -n pinns python=3.8
conda activate pinns
        \end{verbatim}
      \end{block}
    \item Alternativa usando virtualenv:
      \begin{block}{}
        \begin{verbatim}
python3 -m venv pinns_env
source pinns_env/bin/activate
        \end{verbatim}
      \end{block}
  \end{itemize}
\end{frame}

% Frame 2: Instalación de DeepXDE
\begin{frame}{Instalación de DeepXDE}
  \begin{itemize}
    \item DeepXDE: librería principal para PINNs en Python
    \item Instalación vía pip:
      \begin{block}{}
        \begin{verbatim}
pip install deepxde
        \end{verbatim}
      \end{block}
    \item Verificar instalación:
      \begin{block}{}
        \begin{verbatim}
python -c "import deepxde; print(deepxde.__version__)"
        \end{verbatim}
      \end{block}
    \item Dependencias adicionales:
      \begin{itemize}
        \item TensorFlow 2.x (o PyTorch, según backend)
        \item NumPy, SciPy, Matplotlib
      \end{itemize}
  \end{itemize}
\end{frame}

% Frame 3: Estructura básica de un script PINN
\begin{frame}[fragile]{Estructura básica de un script PINN}
  \begin{block}{Ejemplo en Python}
    \begin{verbatim}
import deepxde as dde
import numpy as np

def pde(x, y):
    # Ejemplo: ecuación de Poisson u''(x) = f(x)
    dy_xx = dde.grad.hessian(y, x)
    return dy_xx + np.pi ** 2 * np.sin(np.pi * x)

def boundary(_, on_boundary):
    return on_boundary

def func(x):
    return np.sin(np.pi * x)

    \end{verbatim}
  \end{block}
\end{frame}


\begin{frame}[fragile]{Estructura básica de un script PINN}
  \begin{block}{Ejemplo en Python}
    \begin{verbatim}

# Dominio y condiciones
geom = dde.geometry.Interval(0, 1)
bc = dde.DirichletBC(geom, func, boundary)

# Definición del problema
data = dde.data.PDE(geom, pde, bc, num_domain=50, num_boundary=2)

# Red neuronal
net = dde.maps.FNN([1] + [50] * 3 + [1], "tanh", "Glorot normal")

# Modelo PINN
model = dde.Model(data, net)
model.compile("adam", lr=0.001)

# Entrenamiento
losshistory, _ = model.train(epochs=10000)
    \end{verbatim}
  \end{block}
\end{frame}
